\documentclass{article}
\usepackage{graphicx} 
\usepackage[a4paper, margin=1.2cm]{geometry}

\title{Fase 3. Preparación de los datos}
\author{
    Diana Laura Hernández Almeida \\
    Gonzalo Makenly Higuera Inzunza \\
    José Alfredo López Torres \\
    Juan Pablo De La Peña Gonzalez \\
    Sarah Camila Guzmán Fierro
}
\date{\today}

\begin{document}

\maketitle

\section{Introducción}
El objetivo es construir un conjunto de datos limpio, consistente y adecuado para modelación. Se detallará el proceso ETL (Extract, 
transform, load) durante el reporte, así mismo su justificación y validación en cuanto a conceptos químicos. 

\section{Proceso ETL}
El proceso ETL se compone de tres etapas: Extracción, Transformación y Carga. Su objetivo es limpiar la información, manejar datos
atípicos, imputar valores faltantes y estandarizar las variables para que sean adecuadas para el modelado.
Antes de comenzar con el proceso ETL, se definen reglas basadas en el conocimiento previo. Se define el umbral de correlación en 0.9, 
se definen las variables \texttt{TVOC\_ppb} y \texttt{PM10\_ugm3} para ser eliminadas por su alta correlación, se definen umbrales para cada variable 
basados en la literatura y finalmente se define el umbral de batería baja.

\subsection{Extracción}
El método de extracción carga los archivos \texttt{.xlsx}. Posteriormente separa la columna \texttt{sensor\_periodo} en dos columnas 
independientes: \texttt{sensor} y \texttt{periodo}, y crea un DataFrame por cada sensor solicitado. Si el periodo es el mismo para todos los
registros del sensor, se elimina esa columna. Finalmente ordena los datos cronológicamente usando \texttt{datetime}.

\subsection{Transformación}
La transformación se compone de tres etapas: limpieza, imputación y estandarización/normalización. Esta sigue un patrón de $\texttt{Fit}$ y 
$\texttt{Transform}$, 
es decir, primero aprende los parámetros necesarios y posteriormente los aplica a los datos.

\subsubsection{Limpieza}
Este paso usa los métodos \texttt{\_\_clean\_fit} y \texttt{\_\_clean\_transform} para eliminar datos y decidir qué es un
error de medición.
Primero elimina las columnas altamente correlacionadas. Posteriormente se detectan valores atípicos siguiendo la
siguiente lógica:
Un dato se marca como valor atípico cuando hay discordancia entre el \texttt{AQI} y las concentraciones de contaminantes. Si el \texttt{AQI} es alto, pero no hay
contaminante que lo justifique, se marca como sospechoso. Usando el criterio de la batería se eliminan los datos sospechosos:
Si hay un dato sospechoso y la batería es baja, se clasifica como error de medición y se elimina el dato. Si hay un dato 
sospechoso y la batería es alta, se mantiene el dato. 
\subsubsection{Imputación}
Este paso utiliza los métodos $\texttt{\_\_impute\_fit}$ y $\texttt{\_\_impute\_transform}$.
Después de la limpieza, se verifica que no queden datos faltantes generados por la misma. En caso de que existan, se imputa 
el valor usando interpolación lineal. Este método se elige porque es sencillo, rápido y adecuado para datos temporales. 

\subsubsection{Estandarización y Normalización}
Este paso utiliza los métodos $\texttt{\_\_norm\_std\_fit}$ y $\texttt{\_\_norm\_std\_transform}$.
Para estandarizar y normalizar los datos, estos se separan en dos grupos: contaminantes y variables ambientales/batería.
Para los contaminantes ambientales se utiliza la estandarización \texttt{Z-Score}, la cual resta la media y divide por la
desviación estándar, y se normaliza para que esté entre 0 y 1.
Para la temperatura, humedad y batería se utiliza la normalización Min-Max, la cual resta el valor mínimo y divide por el 
rango. De igual forma, se normaliza para que esté entre 0 y 1.
\subsection{Load}
La última etapa del proceso ETL es la carga, la cual crea una carpeta llamada $\texttt{/Load/}$ si no existe, y exporta cada
DataFrame de sensor procesado a un único archivo $\texttt{.csv}$.
\end{document}